\chapter{Ligament Sprain}

\section*{Definition and Overview}
A ligament sprain is a common musculoskeletal injury characterised by the stretching or tearing of a ligament, which is a tough band of fibrous connective tissue that connects bone to bone and provides structural stability to joints. Sprains typically result from sudden twisting, pulling, or direct trauma that forces a joint beyond its anatomical range of motion. The severity of a sprain is graded from I (mild stretching without joint instability) to II (partial tearing with moderate pain and swelling) and III (complete rupture causing significant instability and loss of function). Sprains commonly affect the ankle, knee, wrist, and thumb. Management ranges from non-invasive conservative measures to surgical reconstruction depending on the location, severity, and functional demands of the patient.

\section*{Epidemiology}
Musculoskeletal injuries, particularly ligament sprains, represent a substantial burden on healthcare systems globally. Ankle sprains are the most common type, accounting for a high percentage of sports-related injuries and emergency department visits. In Australia, millions of sprains occur annually, with a high incidence in active populations, athletes, and manual labourers. Younger individuals, particularly males under 25, have higher rates of sports-related sprains, while older individuals may sustain sprains during daily activities or minor falls. The lateral ligament complex of the ankle, specifically the anterior talofibular ligament (ATFL), is the most frequently injured ligament in the human body.

\section*{Aetiology and Pathophysiology}
Ligament sprains are caused by mechanical forces that exceed the tensile strength of the ligamentous fibers. The common causes and predisposing factors include:
\begin{itemize}
    \item \textbf{Inversion or eversion injuries:} Sudden rolling of the ankle joint, most commonly resulting in inversion which injures the lateral ankle ligaments.
    \item \textbf{Hyperextension or twisting:} Rapid change of direction or pivot movements in the knee, leading to anterior cruciate ligament (ACL) or medial collateral ligament (MCL) sprains.
    \item \textbf{Falls on an outstretched hand (FOOSH):} Impact forces transmitted through the wrist, causing sprains of the scapholunate ligament or wrist capsule.
    \item \textbf{Poor biomechanics and conditioning:} Muscle weakness, poor proprioception, joint laxity, and inadequate warm-up prior to physical activity.
    \item \textbf{Inadequate footwear or uneven surfaces:} Environmental factors that increase the risk of accidental slips, trips, and joint misalignment.
\end{itemize}

\section*{Clinical Features}
The clinical presentation of a ligament sprain varies according to the grade of injury and joint involved. Key features include:
\begin{itemize}
    \item \textbf{Localised pain:} Immediate, sharp pain at the site of injury, which may become a dull ache that worsens with weight-bearing or movement.
    \item \textbf{Swelling and oedema:} Rapid accumulation of fluid around the joint due to acute inflammatory processes.
    \item \textbf{Bruising (ecchymosis):} Discoloration caused by ruptured blood vessels in the damaged ligament or surrounding tissues.
    \item \textbf{Joint instability:} A feeling of the joint "giving way" or being unable to support weight, typical of Grade II and III sprains.
    \item \textbf{Reduced range of motion:} Difficulty moving the joint due to pain, swelling, and structural disruption.
    \item \textbf{Audible pop or tear:} A distinct sound heard or felt at the moment of injury, often indicating a severe or complete ligament rupture.
\end{itemize}

\section*{Differential Diagnosis}
It is critical to differentiate a sprain from other musculoskeletal pathologies:
\begin{itemize}
    \item \textbf{Bone fracture:} Disruption of bone integrity, which must be ruled out using clinical rules (e.g. Ottawa Ankle Rules) or X-rays.
    \item \textbf{Muscle strain:} Injury to muscle fibers or tendons rather than ligaments, often presenting with pain during active muscle contraction.
    \item \textbf{Tendon rupture:} Complete tear of a tendon (e.g. Achilles tendon rupture) mimicking a joint sprain but presenting with distinct functional deficits.
    \item \textbf{Joint dislocation or subluxation:} Complete or partial displacement of articulating bones requiring emergency reduction.
    \item \textbf{Osteochondral lesions:} Damage to joint cartilage and underlying bone, which can occur concurrently with severe sprains.
\end{itemize}

\section*{When to Seek Medical Care}
Minor sprains can often be managed at home, but a professional medical evaluation is recommended if there is severe pain, deformity, or inability to bear weight.

\textbf{Call 000 (or local emergency services) for:}
\begin{itemize}
    \item \textbf{Open wound or exposed tissue:} Severe trauma exposing the joint or bones.
    \item \textbf{Loss of sensation or circulation:} Coldness, numbness, tingling, or blue discoloration distal to the injured joint (neurovascular compromise).
    \item \textbf{Signs of systemic shock:} Extreme dizziness, pale skin, rapid breathing, or confusion following a traumatic fall.
\end{itemize}
Other reasons to see a doctor include inability to take four steps immediately after the injury, severe localised bone tenderness, or lack of improvement after several days of home care.

\section*{Diagnosis and Investigations}
Diagnostic evaluation combines clinical examination with imaging modalities to assess injury severity:
\begin{itemize}
    \item \textbf{Clinical examination:} Palpation for tenderness, assessment of swelling and bruising, and special orthopaedic tests (e.g. anterior drawer test for ankle or knee stability).
    \item \textbf{Plain radiography (X-ray):} Obtained primarily to exclude fractures or joint dislocations, guided by clinical decision rules.
    \item \textbf{Magnetic resonance imaging (MRI):} The gold standard for visualising soft tissues; used to confirm complete tears, assess associated cartilage damage, or plan surgery.
    \item \textbf{Musculoskeletal ultrasound:} A dynamic, non-invasive imaging tool to evaluate ligament integrity and detect joint effusion.
\end{itemize}

\section*{Management}
Management depends on the severity of the sprain and involves conservative or surgical approaches:
\begin{itemize}
    \item \textbf{Conservative management (PRICE/PEACE \& LOVE):} Protect the joint, Elevate above heart level, Avoid anti-inflammatory modalities initially, Compress with a bandage, and Educate the patient. Subsequently, Load the joint progressively and perform Optimistic exercises.
    \item \textbf{Pharmacological management:} Paracetamol is preferred for initial pain relief. NSAIDs may be used after 48 hours to manage persistent swelling, though early use should be avoided to prevent interference with natural tissue healing.
    \item \textbf{Rehabilitation and physiotherapy:} Structured exercises targeting range of motion, muscle strengthening, and proprioception (balance) are vital to restore function and prevent chronic instability.
    \item \textbf{Surgical intervention:} Indicated for certain Grade III sprains (e.g. ACL rupture) in active individuals or when chronic joint instability persists despite rehabilitation.
\end{itemize}

\section*{Complications}
Potential complications of ligament sprains include:
\begin{itemize}
    \item \textbf{Chronic joint instability:} Recurrent sprains and "giving way" due to permanent stretching of ligaments or poor proprioception.
    \item \textbf{Post-traumatic osteoarthritis:} Accelerated joint wear and tear resulting from chronic instability or cartilage damage sustained during the initial injury.
    \item \textbf{Complex regional pain syndrome (CRPS):} A rare, chronic pain condition presenting with severe pain, swelling, and skin changes.
    \item \textbf{Arthrofibrosis:} Joint stiffness and restricted motion caused by excessive scar tissue formation during healing.
    \item \textbf{Persistent pain and swelling:} Chronic inflammation that fails to resolve, often requiring further investigation.
\end{itemize}

\section*{Prognosis and Outcomes}
The prognosis for ligament sprains is generally good, particularly with early and appropriate management. Grade I sprains typically resolve within 1 to 3 weeks. Grade II sprains may require 4 to 8 weeks of rehabilitation. Grade III sprains are more complex; they can take 3 to 6 months or longer to heal, especially if surgical reconstruction is performed. Adherence to a structured rehabilitation programme is the single most important factor in achieving full functional recovery and reducing the risk of re-injury.

\section*{Prevention and Self-Care}
Prevention and active self-care strategies can greatly reduce the risk of sustaining a sprain or experiencing a recurrence:
\begin{itemize}
    \item \textbf{Warm-up and stretch:} Engage in dynamic warm-up exercises before physical activity to prepare muscles and joints.
    \item \textbf{Strengthening and balance training:} Perform targeted exercises to strengthen muscles surrounding vulnerable joints and improve proprioception (e.g. using a wobble board).
    \item \textbf{Appropriate footwear:} Wear supportive, properly fitting shoes designed for the specific activity or sport.
    \item \textbf{Taping or bracing:} Use supportive tape or braces on previously injured joints during high-risk activities to provide mechanical stability.
    \item \textbf{Environmental awareness:} Avoid running or playing sports on uneven, slippery, or poorly lit surfaces.
\end{itemize}

\section*{Sources and Further Reading}
The following reputable organisations provide further information on ligament sprains and joint injuries:
\begin{itemize}
    \item Sports Medicine Australia (SMA). \textit{Injury brochure: ankle sprain and knee sprain.} https://sma.org.au/
    \item Healthdirect Australia. \textit{Sprains and strains overview.} https://www.healthdirect.gov.au/
    \item Better Health Channel. \textit{Sprains and strains: diagnosis, treatment and recovery.} https://www.betterhealth.vic.gov.au/
    \item Mayo Clinic. \textit{Sprains: symptoms, causes, and treatment.} https://www.mayoclinic.org/
    \item NHS. \textit{Sprains and strains information page.} https://www.nhs.uk/
\end{itemize}
